% GENERATED FILE. Edit canonical lessons, then run npm run generate:books.
% canonical-source-hash: fnv1a64:f8dd8261dae9dd5c
% canonical-lessons: GU-C540-baljabri, GU-C540-bandhan, GU-C540-bandhavu, GU-C540-jakadvu, GU-C540-sogand, GU-R540-review-533-540

\chapter{Free or Bound}
\label{ch:gu-free-or-bound}
\hlchaptermodality{540}

\begin{chapteropening}
\textbf{By the end of this chapter:} \emph{I can talk about being forced or bound by a promise.}

Complete the last lesson of chapter 540: I can talk about being forced or bound by a promise.
\end{chapteropening}

\section[\texorpdfstring{\gu{બળજબરી} (baḷjabrī)}{બળજબરી (baḷjabrī)}]{\gu{બળજબરી} (baḷjabrī) — coercion}

\label{lesson:GU-C540-baljabri}



\begin{quote}
Before the new one: say the Gujarati for to take the plunge, then the Gujarati for resolve.
\end{quote}

\subsection*{You'll want to know: \gu{બળજબરી}}
\textbf{\gu{બળજબરી}} — \emph{baḷjabrī} — \textquotedblleft{}coercion\textquotedblright{}.

\begin{quote}
\gu{જો} \gu{તમે} \gu{બળજબરી} \gu{કરશો} \gu{તો} \gu{તે} \gu{કામ} \gu{નહીં} \gu{કરે}.
\end{quote}

\begin{quote}
\emph{jo tame baḷjabrī karśo to te kām nahī̃ kare.}
\end{quote}

\begin{quote}
If you use force, he won't do the work.
\end{quote}

A feminine noun.

\begin{cousinweb}
\textbf{\gu{બળજબરી}} (\emph{baḷjabrī}) is built from \emph{baḷ}, \textquotedblleft{}strength\textquotedblright{}, from Sanskrit \emph{bala} (Monier-Williams), and \emph{jabrī}, usually traced to Arabic \emph{jabr}, \textquotedblleft{}compulsion\textquotedblright{}, through Persian.
\end{cousinweb}

\begin{grammarlens}[title={What you've built}]
One more word for saying what happens if.
\end{grammarlens}

\subsection*{Your turn}
\begin{itemize}
\raggedright
  \item \emph{Say it:} \emph{baḷjabrī}
  \item \emph{Say it:} \emph{baḷjabrī}, once more
  \item \emph{Say it:} \emph{jo tame baḷjabrī karśo to te kām nahī̃ kare.} — If you use force, he won't do the work.
  \item \emph{Recall:} say the Gujarati for to take the plunge, then the Gujarati for resolve, then say \emph{baḷjabrī} again
\end{itemize}

\begin{tcolorbox}[breakable,colback=teal!4,colframe=teal!35!black,title={Before you move on}]
What does \gu{બળજબરી} mean? (\textquotedblleft{}Coercion\textquotedblright{}.) Say it once more.
\end{tcolorbox}
\section[\texorpdfstring{\gu{બંધન} (bandhan)}{બંધન (bandhan)}]{\gu{બંધન} (bandhan) — a tie (that binds)}

\label{lesson:GU-C540-bandhan}



\begin{quote}
Before the new one: say the Gujarati for resolve, then the Gujarati for coercion.
\end{quote}

\subsection*{You'll want to know: \gu{બંધન}}
\textbf{\gu{બંધન}} — \emph{bandhan} — \textquotedblleft{}a tie (that binds)\textquotedblright{}.

\begin{quote}
\gu{જો} \gu{નોકરીનું} \gu{બંધન} \gu{ન} \gu{હોત} \gu{તો} \gu{હું} \gu{આખું} \gu{વરસ} \gu{પ્રવાસ} \gu{કરત}.
\end{quote}

\begin{quote}
\emph{jo nokrīnũ bandhan na hot to hũ ākhũ varas pravās karat.}
\end{quote}

\begin{quote}
If my job didn't tie me down, I would travel all year.
\end{quote}

A neuter noun.

\begin{cousinweb}
\textbf{\gu{બંધન}} (\emph{bandhan}) is a tatsama, taken from Sanskrit \emph{bandhana}, \textquotedblleft{}binding, a bond\textquotedblright{} (Monier-Williams), from the root of \textbf{\gu{બાંધવું}}.
\end{cousinweb}

\begin{grammarlens}[title={What you've built}]
One more word for saying what happens if.
\end{grammarlens}

\subsection*{Your turn}
\begin{itemize}
\raggedright
  \item \emph{Say it:} \emph{bandhan}
  \item \emph{Say it:} \emph{bandhan}, once more
  \item \emph{Say it:} \emph{jo nokrīnũ bandhan na hot to hũ ākhũ varas pravās karat.} — If my job didn't tie me down, I would travel all year.
  \item \emph{Recall:} say the Gujarati for resolve, then the Gujarati for coercion, then say \emph{bandhan} again
\end{itemize}

\begin{tcolorbox}[breakable,colback=teal!4,colframe=teal!35!black,title={Before you move on}]
What does \gu{બંધન} mean? (\textquotedblleft{}A tie (that binds)\textquotedblright{}.) Say it once more.
\end{tcolorbox}
\section[\texorpdfstring{\gu{બંધાવું} (bandhāvũ)}{બંધાવું (bandhāvũ)}]{\gu{બંધાવું} (bandhāvũ) — to be bound (by a promise)}

\label{lesson:GU-C540-bandhavu}



\begin{quote}
Before the new one: say the Gujarati for coercion, then the Gujarati for a tie (that binds).
\end{quote}

\subsection*{You'll want to know: \gu{બંધાવું}}
\textbf{\gu{બંધાવું}} — \emph{bandhāvũ} — \textquotedblleft{}to be bound (by a promise)\textquotedblright{}.

\begin{quote}
\gu{જો} \gu{હું} \gu{વચન} \gu{ન} \gu{આપત} \gu{તો} \gu{હું} \gu{ન} \gu{બંધાત}.
\end{quote}

\begin{quote}
\emph{jo hũ vacan na āpat to hũ na bandhāt.}
\end{quote}

\begin{quote}
If I had not given my word, I would not have been bound.
\end{quote}

Past (he …): \gu{તે} \gu{વચનથી} \gu{બંધાયો}, \emph{te vacanthī bandhāyo}.

\begin{cousinweb}
\textbf{\gu{બંધાવું}} (\emph{bandhāvũ}) is the intransitive partner of \textbf{\gu{બાંધવું}}, \textquotedblleft{}to tie\textquotedblright{}, which CDIAL files under Sanskrit \emph{bandhati}, \textquotedblleft{}binds\textquotedblright{} (CDIAL 9139): to be tied, by a rope or by a promise.
\end{cousinweb}

\begin{grammarlens}[title={What you've built}]
One more word for saying what happens if.
\end{grammarlens}

\subsection*{Your turn}
\begin{itemize}
\raggedright
  \item \emph{Say it:} \emph{bandhāvũ}
  \item \emph{Say it:} \emph{bandhāvũ}, once more
  \item \emph{Say it:} \emph{jo hũ vacan na āpat to hũ na bandhāt.} — If I had not given my word, I would not have been bound.
  \item \emph{Recall:} say the Gujarati for coercion, then the Gujarati for a tie (that binds), then say \emph{bandhāvũ} again
\end{itemize}

\begin{tcolorbox}[breakable,colback=teal!4,colframe=teal!35!black,title={Before you move on}]
What does \gu{બંધાવું} mean? (\textquotedblleft{}To be bound (by a promise)\textquotedblright{}.) Say it once more.
\end{tcolorbox}
\section[\texorpdfstring{\gu{જકડવું} (jakaḍvũ)}{જકડવું (jakaḍvũ)}]{\gu{જકડવું} (jakaḍvũ) — to hold fast}

\label{lesson:GU-C540-jakadvu}



\begin{quote}
Before the new one: say the Gujarati for a tie (that binds), then the Gujarati for to be bound.
\end{quote}

\subsection*{You'll want to know: \gu{જકડવું}}
\textbf{\gu{જકડવું}} — \emph{jakaḍvũ} — \textquotedblleft{}to hold fast\textquotedblright{}.

\begin{quote}
\gu{જો} \gu{તમે} \gu{બાળકનો} \gu{હાથ} \gu{જકડી} \gu{રાખશો} \gu{તો} \gu{તે} \gu{નહીં} \gu{ખોવાય}.
\end{quote}

\begin{quote}
\emph{jo tame bāḷakno hāth jakḍī rākhśo to te nahī̃ khovāy.}
\end{quote}

\begin{quote}
If you hold the child's hand tight, he won't get lost.
\end{quote}

Past, with \gu{તેણે} (he or she …): \gu{તેણે} \gu{હાથ} \gu{જકડ્યો}, \emph{teṇe hāth jakaḍyo}.

\begin{cousinweb}
\textbf{\gu{જકડવું}} (\emph{jakaḍvũ}) is filed by CDIAL under Sanskrit \emph{yatá}, \textquotedblleft{}held, restrained\textquotedblright{}, and CDIAL lists the Gujarati verb, \textquotedblleft{}to tighten, bind\textquotedblright{} (CDIAL 10400).
\end{cousinweb}

\begin{grammarlens}[title={What you've built}]
One more word for saying what happens if.
\end{grammarlens}

\subsection*{Your turn}
\begin{itemize}
\raggedright
  \item \emph{Say it:} \emph{jakaḍvũ}
  \item \emph{Say it:} \emph{jakaḍvũ}, once more
  \item \emph{Say it:} \emph{jo tame bāḷakno hāth jakḍī rākhśo to te nahī̃ khovāy.} — If you hold the child's hand tight, he won't get lost.
  \item \emph{Recall:} say the Gujarati for a tie (that binds), then the Gujarati for to be bound, then say \emph{jakaḍvũ} again
\end{itemize}

\begin{tcolorbox}[breakable,colback=teal!4,colframe=teal!35!black,title={Before you move on}]
What does \gu{જકડવું} mean? (\textquotedblleft{}To hold fast\textquotedblright{}.) Say it once more.
\end{tcolorbox}
\section[\texorpdfstring{\gu{સોગંદ} (sogand)}{સોગંદ (sogand)}]{\gu{સોગંદ} (sogand) — an oath}

\label{lesson:GU-C540-sogand}



\begin{quote}
Before the new one: say the Gujarati for to be bound, then the Gujarati for to hold fast.
\end{quote}

\subsection*{You'll want to know: \gu{સોગંદ}}
\textbf{\gu{સોગંદ}} — \emph{sogand} — \textquotedblleft{}an oath\textquotedblright{}.

\begin{quote}
\gu{જો} \gu{તમે} \gu{સોગંદ} \gu{ખાશો} \gu{તો} \gu{જ} \gu{હું} \gu{માનીશ}.
\end{quote}

\begin{quote}
\emph{jo tame sogand khāśo to ja hũ mānīś.}
\end{quote}

\begin{quote}
Only if you swear will I believe you.
\end{quote}

A masculine noun, used in the plural.

\begin{cousinweb}
\textbf{\gu{સોગંદ}} (\emph{sogand}) is usually traced to Persian \emph{saugand}, \textquotedblleft{}an oath\textquotedblright{}. Gujarati \textquotedblleft{}eats\textquotedblright{} an oath: \textbf{\gu{સોગંદ} \gu{ખાવા}}.
\end{cousinweb}

\begin{grammarlens}[title={What you've built}]
One more word for saying what happens if.
\end{grammarlens}

\subsection*{Your turn}
\begin{itemize}
\raggedright
  \item \emph{Say it:} \emph{sogand}
  \item \emph{Say it:} \emph{sogand}, once more
  \item \emph{Say it:} \emph{jo tame sogand khāśo to ja hũ mānīś.} — Only if you swear will I believe you.
  \item \emph{Recall:} say the Gujarati for to be bound, then the Gujarati for to hold fast, then say \emph{sogand} again
\end{itemize}

\begin{tcolorbox}[breakable,colback=teal!4,colframe=teal!35!black,title={Before you move on}]
What does \gu{સોગંદ} mean? (\textquotedblleft{}An oath\textquotedblright{}.) Say it once more.
\end{tcolorbox}
\section[(dialogue)]{Looking back: from A Roll of the Dice to Free or Bound}

\label{lesson:GU-R540-review-533-540}



\begin{quote}
Say the Gujarati for a die and chaupat.
\end{quote}

\subsection*{Your turn}
Say these aloud:

\begin{itemize}
\raggedright
  \item a die, chaupat, a turn (in a game), a round (of a game), to play a game of chance
  \item gambling, a speculative bet, a wager, a lottery, a chance meeting
  \item an omen, a superstition, a hint, a planet, to make out
  \item an auspicious moment, a choghadiyu, a Hindu almanac, a birth chart, the future
  \item a choice, instead of, favourite, to pick out, to appeal
  \item a quandary, undecided, inner turmoil, to mull over, a settlement
  \item resolute, to settle on, a hard test, to take the plunge, resolve
  \item coercion, a tie (that binds), to be bound, to hold fast, an oath
\end{itemize}

\begin{tcolorbox}[breakable,colback=teal!4,colframe=teal!35!black,title={Before you move on}]
A die? (\textbf{\gu{પાસો}}, \emph{pāso}.) Chaupat? (\textbf{\gu{ચોપાટ}}, \emph{copāṭ}.) A turn (in a game)? (\textbf{\gu{દાવ}}, \emph{dāv}.) A round (of a game)? (\textbf{\gu{બાજી}}, \emph{bājī}.) To play a game of chance? (\textbf{\gu{ખેલવું}}, \emph{khelvũ}.) Gambling? (\textbf{\gu{જુગાર}}, \emph{jugār}.) A speculative bet? (\textbf{\gu{સટ્ટો}}, \emph{saṭṭo}.) A wager? (\textbf{\gu{હોડ}}, \emph{hoḍ}.) A lottery? (\textbf{\gu{લોટરી}}, \emph{loṭrī}.) A chance meeting? (\textbf{\gu{જોગાનુજોગ}}, \emph{jogānujog}.) An omen? (\textbf{\gu{શુકન}}, \emph{śukan}.) A superstition? (\textbf{\gu{વહેમ}}, \emph{vahem}.) A hint? (\textbf{\gu{સંકેત}}, \emph{sanket}.) A planet? (\textbf{\gu{ગ્રહ}}, \emph{grah}.) To make out? (\textbf{\gu{કળવું}}, \emph{kaḷvũ}.) An auspicious moment? (\textbf{\gu{મુહૂર્ત}}, \emph{muhūrt}.) A choghadiyu? (\textbf{\gu{ચોઘડિયું}}, \emph{coghḍiyũ}.) A Hindu almanac? (\textbf{\gu{પંચાંગ}}, \emph{pancā̃g}.) A birth chart? (\textbf{\gu{કુંડળી}}, \emph{kunḍḷī}.) The future? (\textbf{\gu{ભાવિ}}, \emph{bhāvi}.) A choice? (\textbf{\gu{પસંદગી}}, \emph{pasandgī}.) Instead of? (\textbf{\gu{બદલે}}, \emph{badle}.) Favourite? (\textbf{\gu{મનપસંદ}}, \emph{manapsand}.) To pick out? (\textbf{\gu{વીણવું}}, \emph{vīṇvũ}.) To appeal? (\textbf{\gu{રુચવું}}, \emph{rucvũ}.) A quandary? (\textbf{\gu{વિમાસણ}}, \emph{vimāsaṇ}.) Undecided? (\textbf{\gu{ઢચુપચુ}}, \emph{ḍhacupcu}.) Inner turmoil? (\textbf{\gu{ગડમથલ}}, \emph{gaḍamthal}.) To mull over? (\textbf{\gu{ઘોળવું}}, \emph{ghoḷvũ}.) A settlement? (\textbf{\gu{નિવેડો}}, \emph{niveḍo}.) Resolute? (\textbf{\gu{મક્કમ}}, \emph{makkam}.) To settle on? (\textbf{\gu{ઠરાવવું}}, \emph{ṭharāvvũ}.) A hard test? (\textbf{\gu{કસોટી}}, \emph{kasoṭī}.) To take the plunge? (\textbf{\gu{ઝંપલાવવું}}, \emph{jhamplāvvũ}.) Resolve? (\textbf{\gu{નિશ્ચય}}, \emph{niścay}.) Coercion? (\textbf{\gu{બળજબરી}}, \emph{baḷjabrī}.) A tie (that binds)? (\textbf{\gu{બંધન}}, \emph{bandhan}.) To be bound? (\textbf{\gu{બંધાવું}}, \emph{bandhāvũ}.) To hold fast? (\textbf{\gu{જકડવું}}, \emph{jakaḍvũ}.) An oath? (\textbf{\gu{સોગંદ}}, \emph{sogand}.)
\end{tcolorbox}
